\chapter{带外频谱引导的 RF-HSTU 鲁棒建模}
\label{ch:ch3}

本章对应 Paper 1 / IoTJ 主体，关注 \textbf{same-receiver cross-day} 闭集识别。

\section{问题定义}

训练与测试使用同一接收机，测试日（Day5）与训练日不同。评价文件级准确率（File-Acc）与 Macro-F1；24 类 chance level 约 4.17\%。

\section{OOB-guided RF-HSTU 模型}

最终模型 \texttt{F\_cross\_attn\_chirp\_plain} 包含：CNN stem（IQ 局部特征）、RF-HSTU 时序建模、OOB 分支、cross-attention 融合与 chirp-aware 嵌入。图~\ref{fig:ch3_arch} 展示总体架构。

\begin{figure}[htbp]
  \centering
  \includegraphics[width=0.95\linewidth]{fig1_model_architecture.pdf}
  \caption{OOB-guided RF-HSTU 总体架构}
  \label{fig:ch3_arch}
\end{figure}

\textbf{Cross-attention 融合}：in-band token 为 Query，OOB token 为 Key/Value，避免 concat 导致的表征塌缩。测试采用 file-level mean-logits voting。

\section{训练与评估协议}

Manifest：\texttt{cross\_day\_day1to5\_source\_only.csv}；OOB zscore 归一化；batch 128，lr $3\times10^{-3}$，80 epochs，5 seeds。

\section{跨日期主结果}

表~\ref{tab:ch3_main} 与图~\ref{fig:ch3_seeds} 汇总主结果。Ours 达 $75.0\pm5.3\%$，较 CNN-IQ $54.2\pm14.2\%$ 提升约 \textbf{20.8 pp}；bootstrap 95\% CI 约为 $[+9.2,+32.5]$ pp。

\begin{table}[htbp]
  \centering
  \caption{跨日期主结果（Day5 test，mean$\pm$std，\%）}
  \label{tab:ch3_main}
  \begin{tabular}{lcc}
    \toprule
    模型 & File-Acc & File-F1 \\
    \midrule
    CNN-IQ & $54.2\pm14.2$ & $45.6\pm14.8$ \\
    RF-HSTU（无 OOB） & $66.7\pm3.4$ & $59.1\pm4.5$ \\
    \textbf{Ours} & $\mathbf{75.0\pm5.3}$ & $\mathbf{67.9\pm6.8}$ \\
    \bottomrule
  \end{tabular}
\end{table}


\begin{figure}[htbp]
  \centering
  \includegraphics[width=0.85\linewidth]{fig2_cross_day_seed_bars.pdf}
  \caption{跨日期各 seed 文件级准确率}
  \label{fig:ch3_seeds}
\end{figure}

\section{消融实验}

图~\ref{fig:ch3_ablation}：cross-attention 稳定约 75\%；concat fusion collapse 至 9--18\%；chirp 对 cross-attn 均值影响有限。无 OOB 的 RF-HSTU 为 $66.7\pm3.4\%$。

\begin{figure}[htbp]
  \centering
  \includegraphics[width=0.85\linewidth]{fig3_fusion_chirp_ablation.pdf}
  \caption{Fusion 与 Chirp 消融}
  \label{fig:ch3_ablation}
\end{figure}

\section{部署偏移与跨接收机 stress}

图~\ref{fig:ch3_distance} 为距离 LOCO 扩展；图~\ref{fig:ch3_crossrx} 为 strict cross-receiver stress：RX1$\rightarrow$RX2 时 Ours $18.1\pm3.9\%$ vs CNN $4.2\%$，仍远低于 cross-day，且 RX2$\rightarrow$RX1 方向不对称——\textbf{不能声称 receiver-invariant}。

\begin{figure}[htbp]
  \centering
  \includegraphics[width=0.75\linewidth]{fig4_distance_shift.pdf}
  \caption{部署距离偏移（LOCO）}
  \label{fig:ch3_distance}
\end{figure}

\begin{figure}[htbp]
  \centering
  \includegraphics[width=0.75\linewidth]{fig5_cross_receiver_stress.pdf}
  \caption{跨接收机 stress test（limitation）}
  \label{fig:ch3_crossrx}
\end{figure}

\section{本章小结}

OOB-guided cross-attentive RF-HSTU 显著提升同接收机跨日鲁棒性，但 strict cross-receiver 仍 near chance，引出第\ref{ch:ch4}章诊断与校准。
